\documentclass[aspectratio=169]{beamer}
\usepackage[orientation=landscape,size=custom,width=16,height=9,scale=0.5,debug]{beamerposter}
\usepackage[T2A]{fontenc}
\usepackage[utf8]{inputenc}
\usepackage[english,russian]{babel}
\usepackage{cite,enumerate,float,indentfirst}
\usepackage{graphicx}
%\usetheme{umbc2}

\title{Типовые задачи в разработке ПО и их решения}
\author{Alexey Shrub}
\date{2012-09-13}
\begin{document}

\maketitle

\begin{frame}{Введение}
\begin{itemize}
\item Зачем?
\pause
\item Уровни образования
\end{itemize}
\end{frame}

\begin{frame}{Методологии разработки}
\begin{itemize}
\item ООП
\pause
\item Функциональный стиль
\pause
\item Программирование квантовых компьютеров? (Quipper)
\end{itemize}
\end{frame}

\begin{frame}{Coding conventions}
\begin{itemize}
\item Код для человека
\pause
\item Любой дурак может написать код, который понятен компьютеру. Хорошие программисты пишут код, который понятен человеку (Мартин Фаулер)
\pause
\item Именование, отступы
\pause
\item Всегда пишите код так, будто сопровождать его будет склонный к насилию психопат, который знает, где вы живете. — Martin Golding
\end{itemize}
\end{frame}

\begin{frame}{Конфигурирование}
\begin{itemize}
\item Зачем?
\pause
\item Файлы или база?
\pause
\item Форматы файлов и библиотеки.
\end{itemize}
\end{frame}

\begin{frame}{Журналирование}
\begin{itemize}
\item Зачем?
\pause
\item Что? (Уровни)
\pause
\item Куда? (БД или файлы)
\end{itemize}
\end{frame}

\begin{frame}{Профилирование}
\begin{itemize}
\item Зачем?
\pause
\item Инструменты.
\end{itemize}
\end{frame}

\begin{frame}{Протоколы интеграции и маршалинг}
\begin{itemize}
\item Интеграция: файлы, базы, очереди, вызов процедур
\pause
\item Вызовы процедур: SOAP, XML::RPC, JSON::RPC, в какой-то мере REST
\pause
\item Сериализация/Маршалинг: JSON, YAML
\end{itemize}
\end{frame}

\begin{frame}{Хранилища данных}
\begin{itemize}
\item РСУБД (PostgreSQL)
\pause
\item Ключ-значение (Redis, Memcached, Riak)
\pause
\item Документоориентированные (mongoDB)
\pause
\item Очереди
\end{itemize}
\end{frame}

\begin{frame}{Технологии параллелизма}
\begin{itemize}
\item process
\pause
\item thread
\pause
\item event loop
\end{itemize}
\end{frame}

\begin{frame}{Масштабирование ПО}
\begin{itemize}
\item Один в поле не воин
\pause
\item Как распределить запросы (с сессиями и без них)
\pause
\item Репликация. Партицирование
\end{itemize}
\end{frame}

\begin{frame}{Простота и сложность}
\begin{itemize}
\item Отладка кода вдвое сложнее, чем его написание. Так что если вы пишете код настолько умно, насколько можете, то вы по определению недостаточно сообразительны, чтобы его отлаживать. — Brian W. Kernighan.
\pause
\item Есть два способа спроектировать систему. Один — сделать ее настолько простой, что в ней очевидно не будет недостатков, а второй — сделать ее настолько сложной, что в ней не будет очевидных недостатков. Первый способ намного более труден. — Тони Хоар
\pause
\item Unix-way
\end{itemize}
\end{frame}

\begin{frame}{Сложность алгоритмов}
\begin{itemize}
\item Временная сложность алгоритма (в худшем случае)
\begin{itemize}
\item $O(log \: n)$ - бинарный поиск по отсортированному массиву
\item $O(n)$ - поиск в одномерном массиве полным перебором
\item $O(n^2)$ - поиск в двумерном массиве полным перебором, многие виды сортировок
\item $O(c^n)$ - решение задачи коммивояжёра методами динамического программирования
\item $O(n!)$ - решение задачи коммивояжёра полным перебором
\end{itemize}
\pause
\item В лучшем и средних случаях
\pause
\item Пространственная сложность - по раходу памяти
\end{itemize}
\end{frame}

\begin{frame}{Оптимизация}
\begin{itemize}
\item Всему своё время
\pause
\item Компилятор умный
\pause
\item Узкие места на нужном языке
\end{itemize}
\end{frame}

\begin{frame}{Установка ПО}
\begin{itemize}
\item Версии
\pause
\item Зависимости
\pause
\item Откат
\end{itemize}
\end{frame}

\begin{frame}{Локализация и интернационализация (технологии и подходы)}
\begin{itemize}
\item Интернационализация - чтоб могло (utf8, gettext, .po файлы и инструменты для них)
\pause
\item Локализация - для конкретной страны (перевод, форматы дат, денег)
\end{itemize}
\end{frame}

\begin{frame}{Конечные автоматы?}
\begin{itemize}
\item Регулярные выражения
\pause
\item PCRE
\end{itemize}
\end{frame}

\begin{frame}{Парсинг}
\begin{itemize}
\item Нужен? DSL?
\pause
\item lex/bison/yacc
\pause
\item бинарные форматы (Kaitai Struct)
\end{itemize}
\end{frame}


\begin{frame}{"Читай много, но не очень много книг." Бенджамин Франклин}
\begin{itemize}
\item "Совершенный код" Стив Макконнелл
\pause
\item "Программист-прагматик. Путь от подмастерья к мастеру" Э. Хант, Д. Томас
\pause
\item "Рефакторинг. Улучшение существующего кода" Мартин Фаулер
\pause
\item "Приемы объектно-ориентированного проектирования. Паттерны проектирования" Э. Гамма, Р. Хелм, Р. Джонсон, Дж. Влиссидес
\pause
\item "Этюды для программистов" Чарльз Уэзерелл
\pause
\item "Искусство программирования" Дональд Кнут
\pause
\item "Регулярные выражения" Джеффри Фридл
\end{itemize}
\end{frame}

\begin{frame}{Что ещё?}
"I belive that you can never be a good programmer if you know only the syntax of the language but not how the compilation or run-time enviroment is implemented. For example, ... a Java programmer should know why a thread may never get control in a uniprocessor setup even if it is not blocked."

"I am of the opinion that an application developer should master at least the following six major areas: user interfaces, persistance, interprocess communication and networking, parsing and code generation, the Web, and the operating system"

Sriram Srinivasan "Advanced Perl Programming"
\end{frame}

\begin{frame}{Вопросы?}
\begin{block}{Исходники презентации (LaTeX, Beamer):}
https://github.com/worldmind/typical-dev-tasks-presentation-ru.git
\begin{center}
\includegraphics{qr-git-url.png}
\end{center}
\end{block}
\begin{block}{Feedback to: ashrub@yandex.ru}
\end{block}

\end{frame}

\end {document}

